% ============================================================
%  Chapter: Literature Review (Condensed)
%  Table retained, verbose prose trimmed
% ============================================================
\chapter{Literature Review}
\label{chap:related_work}

This chapter surveys existing research in serverless analytics, stream processing,
resource allocation, and reinforcement learning for cloud/edge computing. Each section
examines a cluster of related works, and a comparative table at the end systematically
identifies the research gaps motivating the proposed contributions.

% ─────────────────────────────────────────────────────────────
\section{Serverless Computing and Analytics}
% ─────────────────────────────────────────────────────────────
Serverless computing — also known as Function-as-a-Service (FaaS) — abstracts server
provisioning entirely, scaling applications automatically and charging only for actual
compute time~\cite{hendrickson2020serverless}. Commercial serverless analytics offerings
such as AWS EMR Serverless, Google Dataproc Serverless, and Azure Synapse Serverless
Spark Pools provide managed Spark environments but lack intelligent optimisation:
resource allocation follows simple heuristics or user-specified configurations that
cannot adapt dynamically to workload
characteristics~\cite{aws_emr_serverless_docs,dataproc_serverless_docs,azure_synapse_docs}.

Academic works further expose this gap. Pu et al.\ proposed Locus for data-shuffle
optimisation~\cite{pu2019shuffling}; Wawrzoniak et al.\ investigated running
unmodified Spark on serverless platforms~\cite{wawrzoniak2024off}; Werner presented a
reference architecture for serverless big-data processing~\cite{werner2024reference}.
These works collectively highlight the disconnect between current FaaS platforms and
large-scale analytics workloads.

% ─────────────────────────────────────────────────────────────
\section{Shuffle Optimisation in Distributed Analytics}
% ─────────────────────────────────────────────────────────────
Shuffle operations (data redistribution during \texttt{groupBy}/\texttt{join}) are a
major bottleneck in distributed analytics. Eizaguirre \&
Sánchez-Artigas~\cite{eizaguirre2022seer} presented \textbf{Seer}, an auto-tuned
shuffle mechanism that dynamically selects between in-memory and object-storage
back-ends, achieving 23--45\% improvements over serverful Spark on TPC-DS benchmarks.
However, Seer is limited to a static two-tier storage model without intermediate tiers
(e.g.\ NVMe), lacks workload-specific learning, and does not integrate with resource
allocation optimisation. The proposed framework extends this by incorporating four
storage tiers (Redis, NVMe, S3, Glacier) with LSTM-based temperature
prediction~\cite{lstm_storage_temp2022}.

Kim et al.\ introduced Flint, a prototype Spark engine on AWS
Lambda~\cite{kim2021function}, and Wawrzoniak et al.~\cite{wawrzoniak2024off} explored
running unmodified Spark/Drill on Lambda. These works underline the need for specialised
adaptation when scaling analytics engines onto serverless infrastructure.

% ─────────────────────────────────────────────────────────────
\section{Resource Allocation and Auto-scaling}
% ─────────────────────────────────────────────────────────────
Nestorov et al.~\cite{nestorov2024dexter} presented \textbf{Dexter}, a resource
allocator for serverless analytics achieving up to 4.65$\times$ cost reduction and
3.50$\times$ performance-cost efficiency improvement via predictive linear-regression
and decision-tree models. However, Dexter's linear models cannot capture complex
non-linear relationships, it optimises cost alone without latency SLA or throughput
considerations, and its reactive approach requires workload execution before adaptation.
The proposed work addresses these limitations through deep RL (PPO) for simultaneous
multi-objective optimisation~\cite{schulman2017ppo}.

In the broader DRL space, Cui et al.~\cite{cui2023drl} apply DRL to cloud--edge
collaborative resource allocation in IoV systems; Fu et al.~\cite{fu2022distributed}
analyse memory allocation using RL for edge-PLC resources; Zhou et
al.~\cite{zhou2024deep} survey DRL-based cloud scheduling methods. These works
demonstrate the maturity of RL for cloud/edge resource management but few target
serverless Spark environments specifically.

% ─────────────────────────────────────────────────────────────
\section{Stream Processing Frameworks}
% ─────────────────────────────────────────────────────────────
Cai et al.~\cite{cai2024spsc} presented \textbf{SPSC}, a stream-processing framework
on AWS Lambda achieving $\sim$10\% improvement over Alibaba Flink with automatic
parallelisation and true serverless deployment. However, Lambda's 15-minute execution
limit, stateless functions requiring external state-stores, and limited support for
complex windowed operations make it unsuitable for the rich stateful analytics required
by IoT workloads. The proposed framework uses serverless Spark instead, enabling
SQL-like queries and complex stateful operations while maintaining serverless benefits.

Rotchford et al.~\cite{rotchford2024laminar} introduced \textbf{Laminar 2.0}, improving
developer productivity through code recommendations for serverless stream processing.
While valuable for usability, Laminar 2.0 does not address resource optimisation or
cost-efficiency — the primary focus of this work.

% ─────────────────────────────────────────────────────────────
\section{ServerlessYarn: Hybrid Serverless Platform}
% ─────────────────────────────────────────────────────────────
Castellanos-Rodríguez et al.~\cite{castellanos2024serverlessyarn} presented
\textbf{ServerlessYarn}, achieving up to 41\% runtime reduction and 50\% resource
savings through automatic container scaling on YARN clusters. However, it remains
pseudo-serverless (still requires YARN cluster deployment), lacks intelligent
optimisation beyond reactive scaling, and does not support multi-objective trade-offs.
True serverless platforms eliminate all cluster management, and the proposed framework
builds upon these with intelligent RL-based optimisation layers.

% ─────────────────────────────────────────────────────────────
\section{Reinforcement Learning in Cloud Computing}
% ─────────────────────────────────────────────────────────────
RL has emerged as a promising paradigm for cloud optimisation. Liu et
al.~\cite{liu2020edge} apply multi-agent RL for IoT edge network allocation; Zhou et
al.~\cite{zhou2024deep} survey DRL-based cloud scheduling. However, existing efforts
mainly target containerised or VM-based clusters rather than serverless Spark.

Proximal Policy Optimisation (PPO)~\cite{schulman2017ppo} has become a standard
algorithm for continuous-control tasks due to its sample-efficiency and stability. The
proposed work leverages PPO's multi-objective capabilities to simultaneously optimise
cost, latency, and throughput in serverless Spark — a novel application domain.

% ─────────────────────────────────────────────────────────────
\section{Edge Computing and IoT Data Processing}
% ─────────────────────────────────────────────────────────────
Al-Dulaimy et al.~\cite{alDulaimy2024continuum} survey edge--cloud continuum
architectures, and Kjorveziroski et al.~\cite{kjorveziroski2021iot} map progress in
serverless computing at the edge. These works motivate the proposed CRDT-based
synchronisation strategy for IoT-driven analytics pipelines, addressing the gap between
purely cloud-centric frameworks and latency-sensitive edge deployments.

% ─────────────────────────────────────────────────────────────
\section{Research Gaps and Motivation}
% ─────────────────────────────────────────────────────────────
Table~\ref{tab:literature_gaps} summarises the capabilities of existing works and the
gaps addressed by this thesis.

\begin{table}[htbp]
\centering
\caption{Comparative Analysis of Existing Works and Research Gaps}
\label{tab:literature_gaps}
\small
\begin{tabular}{|p{2.5cm}|p{1.5cm}|p{1.5cm}|p{1.5cm}|p{1.5cm}|p{1.5cm}|p{1.6cm}|}
\hline
\textbf{Feature} & \textbf{Seer} & \textbf{Dexter} & \textbf{SPSC} & \textbf{Yarn} & \textbf{Laminar 2.0} & \textbf{Proposed} \\
\hline
\textbf{Serverless Spark} & Yes & Yes & No & No & No & Yes \\
\hline
\textbf{IoT-Specific} & No & No & Partial & No & No & Yes \\
\hline
\textbf{Multi-Objective Optimization} & No & No & No & No & No & Yes \\
\hline
\textbf{RL-Based Allocation} & No & No & No & No & No & Yes \\
\hline
\textbf{Edge-Cloud Sync} & No & No & No & No & No & Yes \\
\hline
\textbf{CRDT Integration} & No & No & No & No & No & Yes \\
\hline
\textbf{Predictive Prefetching} & No & No & No & No & No & Yes \\
\hline
\textbf{Cross-Cloud Support} & No & No & No & No & No & Yes \\
\hline
\textbf{Cost Optimization} & No & Yes & No & Partial & No & Yes \\
\hline
\textbf{Latency Optimization} & Partial & No & No & No & No & Yes \\
\hline
\textbf{Throughput Optimization} & Partial & No & Partial & No & No & Yes \\
\hline
\end{tabular}
\end{table}

\noindent The key research gaps are:
\begin{enumerate}
  \item \textbf{No multi-objective optimisation:} Existing works optimise single
        objectives (Dexter: cost; Seer: shuffle), whereas IoT applications require
        balancing cost, latency, throughput, and reliability simultaneously.
  \item \textbf{No RL in serverless Spark:} While DRL has been applied to cloud/edge
        allocation, no prior work targets serverless Spark streaming with RL.
  \item \textbf{Limited storage tiers:} Seer's two-tier model misses cost-performance
        opportunities available through NVMe and archival tiers.
  \item \textbf{No edge--cloud integration:} Existing frameworks focus on cloud-only
        processing, ignoring edge benefits for latency-sensitive IoT.
  \item \textbf{Vendor lock-in:} Most frameworks target a single cloud provider;
        cross-cloud portability remains unaddressed.
  \item \textbf{No IoT-specific optimisation:} General analytics frameworks do not
        exploit IoT domain knowledge (burst patterns, sensor heterogeneity, temporal
        correlations).
\end{enumerate}

% ─────────────────────────────────────────────────────────────
\section{Summary}
% ─────────────────────────────────────────────────────────────
This chapter surveyed literature across serverless computing, stream processing,
resource allocation, and reinforcement learning. While significant progress exists in
individual areas, no prior work provides an integrated solution covering the full
spectrum of real-time IoT analytics challenges in serverless environments. The proposed
framework — combining RL-based multi-objective optimisation, four-tier adaptive shuffle,
CRDT-based edge--cloud synchronisation, and cross-cloud portability — addresses these
gaps comprehensively.
